% ============================================================
%  config.inc.tex — Настраиваемые параметры шаблона ЕДТ
%  Меняй значения здесь, не трогая preamble.inc.tex / edtstyle.sty
% ============================================================

% --- Нумерация (раскомментировать = нумерация внутри раздела) ---
% \EqInChapter        % формулы:  1.1, 1.2, 2.1 ...
% \TableInChapter     % таблицы:  1.1, 1.2, 2.1 ...
% \PicInChapter       % рисунки:  1.1, 1.2, 2.1 ...
\ListingInChapter     % листинги: 1.1, 1.2, 2.1 ...

% --- Поля страницы (как в Word-оригинале ЕДТ: 25 мм) ---
\geometry{left=25mm}
\geometry{right=25mm}
\geometry{top=25mm}
\geometry{bottom=20mm}
\geometry{headheight=14pt}
\geometry{headsep=8mm}
\geometry{footskip=12mm}

% --- Шрифты ---
\setmainfont{Times New Roman}   % основной текст
\setsansfont{Arial}             % заголовки, колонтитулы, подписи

% --- Кегль основного текста задаётся опцией класса в rpz.tex (12pt) ---

% ============================================================
%  Реквизиты документа (титул + колонтитул)
% ============================================================
\renewcommand{\edtProjectName}{Интеллектуальная система управления оповещениями}
\renewcommand{\edtBannerTitle}{Единый документ требований}
\renewcommand{\edtDocTitle}{Единый документ требований к ИТ-решению}
\renewcommand{\edtDocVersion}{1.1}
\renewcommand{\edtDocDate}{08.08.2026}
\renewcommand{\edtDocPurpose}{Определяет контекст для разработки требований и
описывает требования различных уровней, обеспечивая их согласованность и
прослеживаемость для успешной реализации проекта ИТ и ЦТ}
% Строка верхнего колонтитула (по умолчанию — «проект - ЕДТ - версия X»)
\renewcommand{\edtRunningHead}{\edtProjectName{} - ЕДТ - версия \edtDocVersion}

% ============================================================
%  Текстовые константы (названия разделов, подписи)
% ============================================================

% --- Оглавление ---
\addto\captionsrussian{\renewcommand\contentsname{Содержание}}

% --- Введение ---
\renewcommand\Introduction{\chapter*{Введение}\addcontentsline{toc}{chapter}{Введение}}

% --- Заключение ---
\renewcommand\Conclusion{\chapter*{Заключение}\addcontentsline{toc}{chapter}{Заключение}}

% --- Список источников ---
\addto\captionsrussian{\renewcommand\bibname{Список использованных источников}}

% --- Приложение ---
\addto\captionsrussian{\renewcommand\appendixname{Приложение}}

% --- Подписи к плавающим объектам ---
\addto\captionsrussian{\renewcommand\tablename{Таблица}}
\addto\captionsrussian{\renewcommand\figurename{Рисунок}}
\SetupFloatingEnvironment{listing}{name=Листинг}
